\documentclass[10pt,a4paper]{amsart}
\usepackage[margin=19mm]{geometry}
\usepackage{amsmath,amssymb,mathtools}
\usepackage[scr=rsfs,cal=euler]{mathalfa}
\usepackage{xfrac}
\usepackage[unicode,hidelinks,bookmarks=false]{hyperref}
\newcommand{\ie}{i.e.\ }
\newcommand{\cf}{cf.\ }
\newcommand{\resp}{respectively\ }
\newcommand{\Subsection}{\subsection}
\providecommand{\nobreakdash}{\nobreak\hbox{-}}
\setlength{\emergencystretch}{2em}
\multlinegap0pt
\usepackage{mathtools}
\usepackage{amssymb}
\usepackage[shortlabels]{enumitem}
\usepackage{xcolor}
\newtheorem{theorem}{Theorem}
\title{Spectral Perspectives on FAT Graph Colorings}
\author{Lies Beers, Raffaella Mulas}
\date{Source: arXiv:2604.21665v1 (CC BY 4.0)}
\begin{document}
\maketitle

\noindent\emph{Source and licence.} Spectral Perspectives on FAT Graph Colorings. Version arXiv:2604.21665v1 (CC BY 4.0), distributed by arXiv under the Creative Commons Attribution 4.0 International licence, http://creativecommons.org/licenses/by/4.0/, as recorded in the arXiv licence metadata for this exact version and shown on the abstract page https://arxiv.org/abs/2604.21665v1. Full licence notice: \texttt{licenses/arxiv-2604.21665-CC-BY-4.0.txt}.\par\smallskip

\noindent\textit{Editorial scope.} Counted unit: the article's theorem thm:removing (source0.tex lines 726--731). The article's complete original proof of it is delivered with it (source0.tex lines 732--757, one \texttt{proof} environment of 26 lines). No mathematics is written, repaired or re-derived: the statement and the proof are the article's own lines, in the article's own notation, and no retained line is substituted or re-flowed.  No further source-local context is needed: every cross-reference of every retained line resolves inside the two delivered spans.  Nothing was omitted from any retained span, so the package carries no omission record.  No retained line contains a citation, so the wrapper carries no bibliography and no bibliography entry.  No retained line contains a figure environment, an image-inclusion command or drawing code, so no image is delivered and the package declares no asset dependency.  Editorial formatting only, in addition: an AMS \texttt{amsart} preamble that restates the article's own declarations and macros used by the retained lines, namely the environments \texttt{theorem} on separate counters, together with the standard packages those lines need.  The retained environments print in this document as Theorem 1 in order of appearance, whereas the article prints the same items with the numbering of its own full document.  The complete native source of the article as distributed in the arXiv e-print for this exact version is bundled unchanged as \texttt{sources/199042/source0.tex}.
\begin{theorem}\label{thm:removing}
    Let $V_1,\ldots,V_k$ denote the coloring classes of $G$ with respect to a FAT $k$-coloring with parameter $\alpha$, and let $I\subseteq \{1,\ldots,k\}$ be an index set of size $\lvert I\rvert\leq k-2$. Then, the induced subgraph $G\bigl(V\setminus\bigcup_{i\in I}V_i\bigr)$ admits a FAT $(k-\lvert I\rvert)$-coloring with parameter
    \[
        \frac{\alpha}{1-\alpha\lvert I\rvert}.
    \]
\end{theorem}

\begin{proof}
    Let \[
        G_{I}\coloneqq G\bigl(V\setminus\bigcup_{i\in I}V_i\bigr),
    \] and fix $v\in V(G_I)$.
    In the original graph $G$, the vertex $v$ has $\deg_G(v)$ neighbors. In the induced subgraph $G_I$, all neighbors of $v$ that lie in the coloring classes indexed by $I$ are removed. Since the given coloring is FAT with parameter $\alpha$, the vertex $v$ has exactly $\alpha\,\deg_G(v)$ neighbors in each coloring class $V_i$ for which $v\notin V_i$. In particular, for each $i\in I$, $v$ has $\alpha\,\deg_G(v)$ neighbors in $V_i$, and hence
    \[
        \deg_{G_I}(v)
        \;=\;
        \deg_G(v) - \alpha\lvert I\rvert\,\deg_G(v)
        \;=\;
        \bigl(1-\alpha\lvert I\rvert\bigr)\deg_G(v).
    \]
    Moreover, for any coloring class $V_i$ with $i\notin I$ and $v\notin V_i$, the vertex $v$ still has $\alpha\,\deg_G(v)$ neighbors in $V_i$. Rewriting in terms of $\deg_{G_I}(v)$ yields
    \[
        \alpha\,\deg_G(v)
        \;=\;
        \frac{\alpha}{1-\alpha\lvert I\rvert}\,\bigl(1-\alpha\lvert I\rvert\bigr)\deg_G(v)
        \;=\;
        \frac{\alpha}{1-\alpha\lvert I\rvert}\,\deg_{G_I}(v).
    \]
    Consequently, in $G_I$ the vertex $v$ has exactly
    \[
        \frac{\alpha}{1-\alpha\lvert I\rvert}\,\deg_{G_I}(v)
    \]
    neighbors in each remaining coloring class $V_i$ with $i\notin I$ and $v\notin V_i$. Hence, the sets $V_i$ for $i\in \{1,\ldots,k\}\setminus I$ form a FAT $(k-\lvert I\rvert)$-coloring of $G_I$ with parameter $\alpha/(1-\alpha\lvert I\rvert)$.
\end{proof}

\end{document}
